% SC1008 — Chapter 5: C++ Classes & OOP Basics (0->100 ground-up)
\chapter{C++ Classes and Object-Oriented Basics}
\label{ch:cpp-classes-oop}

\begin{quote}
\emph{``C++ is C with a better abstraction toolkit.''} A \texttt{class} bundles data with the operations that maintain its invariants --- turning a loose struct into a disciplined object.
\end{quote}

\noindent\fbox{\parbox{\dimexpr\linewidth-2\fboxsep-2\fboxrule}{%
\textbf{From zero: C structs $\to$ C++ classes.} We start with why bare structs are error-prone (anyone can break invariants), then add constructors, encapsulation, and methods step by step. No prior OOP assumed.}}

\section*{Learning Outcomes}
After this chapter you will be able to:
\begin{enumerate}[label=(L\arabic*)]
    \item define a C++ class with access control (\texttt{private}/\texttt{public});
    \item write constructors, destructors, and use member-initialiser lists;
    \item distinguish value vs.\ reference semantics and use \texttt{const} correctly;
    \item explain encapsulation and class invariants;
    \item trace object lifetime (construction $\to$ use $\to$ destruction).
\end{enumerate}

% ============================================================
\section{From Struct to Class}
\label{sec:struct-to-class}
A C struct is passive: any code can set \texttt{s.gpa = 999}. A class makes data \texttt{private} and exposes only safe operations.

\begin{figure}[h]\centering
\begin{tikzpicture}[scale=0.70,every node/.style={font=\scriptsize}]
\node[draw,thick,rounded corners=3pt,fill=gray!10,minimum width=3.2cm,minimum height=1.6cm,align=left] (bad) at (0,0) {\textbf{struct} (open)\\ \texttt{gpa = 999} allowed\\ no invariant\\ {\tiny anyone can corrupt}};
\node[draw,thick,rounded corners=3pt,fill=green!12,minimum width=3.2cm,minimum height=1.6cm,align=left] (good) at (5.0,0) {\textbf{class} (guarded)\\ \texttt{private: gpa}\\ \texttt{public: setGPA()}\\ checks $0\le$\texttt{gpa}$\le5$};
\draw[-{Latex},thick,blue!60!black] (bad)--(good) node[midway,above,font=\tiny] {encapsulate};
\draw[thick,red!60!black,dashed] (5.0,-0.85)--(5.0,0.85);
\node[font=\tiny,fill=orange!15,rounded corners=2pt,inner sep=2pt] at (5.0,1.25) {boundary};
\node[font=\tiny,fill=yellow!12,rounded corners=2pt,inner sep=2pt] at (0,-1.25) {no boundary};
\end{tikzpicture}
\caption{Encapsulation: a class erects a wall. Outside code must use the public interface, which can enforce invariants.}
\label{fig:encapsulation}
\end{figure}

\begin{lstlisting}[language=C++]
#include <string>
class Student{
private:
    int id_;
    std::string name_;
    double gpa_;                 // invariant: 0.0 <= gpa_ <= 5.0
public:
    Student(int id, std::string name, double gpa) // constructor
      : id_(id), name_(std::move(name)), gpa_(gpa) {
        if(gpa_<0 || gpa_>5) throw std::out_of_range("gpa");
    }
    double gpa() const { return gpa_; }          // getter, const = won't mutate
    void setGPA(double g){                       // setter with check
        if(g<0||g>5) throw std::out_of_range("gpa");
        gpa_=g;
    }
    std::string name() const { return name_; }
};
\end{lstlisting}
\texttt{class} defaults to \texttt{private}; \texttt{struct} defaults to \texttt{public} --- otherwise identical in C++.

\section{Constructors, Destructors, Initialiser Lists}
\label{sec:ctor-dtor}

\begin{figure}[h]\centering
\begin{tikzpicture}[scale=0.68,every node/.style={font=\scriptsize}]
\node[draw,thick,fill=blue!10,rounded corners=3pt,minimum width=2.0cm,minimum height=0.65cm] (def) at (0,0) {\texttt{Student s;}};
\node[draw,thick,fill=orange!15,rounded corners=3pt,minimum width=2.6cm,minimum height=0.65cm] (par) at (3.8,0) {\texttt{Student s(1,"Asha",4.8)}};
\node[draw,thick,fill=green!12,rounded corners=3pt,minimum width=2.0cm,minimum height=0.65cm] (cpy) at (7.6,0) {\texttt{Student t=s;}};
\node[draw,thick,fill=violet!12,rounded corners=3pt,minimum width=3.0cm,minimum height=0.65cm] (des) at (3.8,-1.35) {destructor \texttt{\~{}Student()} at scope end};
\draw[-{Latex},thick] (def)--(par);
\draw[-{Latex},thick] (par)--(cpy);
\draw[-{Latex},thick,red!60!black] (par)--(des);
\draw[-{Latex},thick,red!60!black] (cpy)--(des);
\node[font=\tiny,fill=yellow!12,rounded corners=2pt,inner sep=2pt] at (3.8,0.85) {construction (initialiser list runs first)};
\node[font=\tiny,fill=yellow!12,rounded corners=2pt,inner sep=2pt] at (3.8,-2.05) {destruction in reverse order of construction};
\end{tikzpicture}
\caption{Object lifetime: construction via initialiser list, then use, then automatic destruction when the scope ends. Copy construction is discussed in Ch.~8.}
\label{fig:ctor-dtor}
\end{figure}

\begin{lstlisting}[language=C++]
class Vec2{
    double x_, y_;
public:
    Vec2(): x_(0), y_(0) {}                 // default constructor
    Vec2(double x,double y): x_(x), y_(y) {} // parameterised
    // member initialiser list : x_(x) runs BEFORE body -- required for const/ref members
    double norm() const { return std::sqrt(x_*x_+y_*y_); }
    ~Vec2(){ /* release resources if any */ } // destructor
};
Vec2 a;              // default
Vec2 b(3,4);         // param
Vec2 c = b;          // copy (memberwise for now; Ch.8 refines)
\end{lstlisting}
Prefer initialiser lists over assignment in the body: they directly construct members, while assignment constructs then overwrites.

\section{Encapsulation and Invariants}
\label{sec:invariants-oop}
A class invariant is a predicate that every public operation must preserve (e.g.\ $0\le$\texttt{gpa}$\le5$, vector \texttt{size}$\le$\texttt{capacity}). By making data private and checking in setters/constructors, the invariant holds for all clients. Think of the class as a tiny machine with a valid-state set; public methods are the only buttons.

\begin{lstlisting}[language=C++]
class BoundedCounter{
    int val_, max_;
public:
    BoundedCounter(int m): val_(0), max_(m) {}
    void inc(){ if(val_<max_) ++val_; }
    void reset(){ val_=0; }
    int get() const { return val_; }
    // invariant: 0 <= val_ <= max_  (always, between calls)
};
\end{lstlisting}

\section{Const Correctness and References}
\label{sec:const-ref}
\begin{lstlisting}[language=C++]
void print(const Student& s){          // const ref: no copy, no mutate
    std::cout << s.name() << " " << s.gpa() << "\n";
}
Student s(1,"Asha",4.8);
print(s);                              // binds without copying string
// s.gpa() is const member function -- callable on const Student&
\end{lstlisting}
\texttt{const} on a method means \texttt{this} is \texttt{const}; it can be called on const objects. Pass large objects by \texttt{const\&} to avoid copies. Return by value for small types, by const ref for members (with lifetime care).

\begin{figure}[h]\centering
\begin{tikzpicture}[scale=0.70,every node/.style={font=\scriptsize}]
\draw[thick,fill=blue!10] (0,0) rectangle (2.4,0.70);
\node at (1.2,0.35) {\texttt{Student s}};
\draw[thick,fill=orange!15] (3.2,0) rectangle (5.6,0.70);
\node at (4.4,0.35) {\texttt{const Student\& r = s}};
\draw[-{Latex},thick,red!70!black] (4.4,0.70)--(1.8,0.70) node[midway,above,font=\tiny] {alias, no copy};
\node[font=\tiny,fill=green!12,rounded corners=2pt,inner sep=2pt] at (1.2,-0.45) {owns data};
\node[font=\tiny,fill=violet!12,rounded corners=2pt,inner sep=2pt] at (4.4,-0.45) {view, cannot mutate};
\end{tikzpicture}
\caption{Reference semantics: \texttt{const Student\&} is an alias (no copy) with read-only access. Value semantics would copy the whole object.}
\label{fig:ref-semantics}
\end{figure}


\section{Operator Overloading}
Classes can give meaning to operators like \texttt{+}, \texttt{==}, \texttt{<<} so objects feel like built-in types.

\begin{lstlisting}[language=C++]
class Vec2{
    double x_,y_;
public:
    Vec2(double x=0,double y=0): x_(x), y_(y) {}
    Vec2 operator+(const Vec2& o) const { return {x_+o.x_, y_+o.y_}; }
    bool operator==(const Vec2& o) const { return x_==o.x_ && y_==o.y_; }
    double x() const { return x_; }
    double y() const { return y_; }
    friend std::ostream& operator<<(std::ostream& os, const Vec2& v){
        return os<<"("<<v.x_<<","<<v.y_<<")";
    }
};
Vec2 a(1,2), b(3,4);
Vec2 c = a+b;            // (4,6)
std::cout<<c<<"\n";       // (4,6)
\end{lstlisting}
Overload only when semantics are obvious; \texttt{+} should not surprise. Stream insertion \texttt{operator<<} is idiomatically a \texttt{friend} non-member.

\section{Static Members and Friends}
\texttt{static} members belong to the class, not to any instance; \texttt{friend} grants selective access to private data.

\begin{lstlisting}[language=C++]
class Counter{
    static int total_;   // shared across all objects
    int id_;
public:
    Counter(): id_(++total_) {}
    static int total(){ return total_; }
    friend void debug(const Counter& c); // can read c.id_ private
};
int Counter::total_ = 0; // define static outside class
void debug(const Counter& c){ std::cout<<c.id_<<"\n"; }
\end{lstlisting}
Static data needs an out-of-class definition (or \texttt{inline static} since C++17). Friends break encapsulation minimally --- prefer member functions when possible.


\section{Member Functions vs.\ Non-Members}
Not every operation should be a member. Prefer non-member non-friend when it can be implemented via the public interface --- it reduces coupling and keeps the class minimal.

\begin{lstlisting}[language=C++]
class Vec2{
    double x_,y_;
public:
    Vec2(double x=0,double y=0): x_(x), y_(y) {}
    double x() const { return x_; }
    double y() const { return y_; }
    Vec2& operator+=(const Vec2& o){ x_+=o.x_; y_+=o.y_; return *this; }
};
// non-member built from public members + +=
Vec2 operator+(Vec2 a, const Vec2& b){ a+=b; return a; }
double dot(const Vec2& a, const Vec2& b){ return a.x()*b.x()+a.y()*b.y(); }
bool isZero(const Vec2& v){ return v.x()==0 && v.y()==0; }
\end{lstlisting}
If a function needs private access, make it a member or friend; otherwise keep it outside. This is the ``minimal interface'' principle (Meyers).

\section{Chapter Notes}
Stroustrup \emph{The C++ Programming Language} (4th ed.) and Meyers \emph{Effective C++} are the classic guides. Rule of Zero/Three/Five (when to write copy/move/destructor) is covered in Ch.~8.

\section{Exercises}
\begin{enumerate}[label=E5.\arabic*, leftmargin=*, itemsep=0.6em]
\item \textbf{Invariant.} Design \texttt{class Date\{int d,m,y\}} with invariant ``valid calendar date''. Which operations must validate? Sketch constructor and \texttt{setDay}.
\item \textbf{Initialiser list.} Why must a \texttt{const int id\_} or \texttt{int\& ref\_} member be initialised in the member-initialiser list, not in the body?
\item \textbf{Const.} Mark which methods should be \texttt{const} in a \texttt{BankAccount} class (\texttt{balance()}, \texttt{deposit()}, \texttt{withdraw()}, \texttt{getId()}). What breaks if you omit \texttt{const}?
\item \textbf{Encapsulation break.} Show how returning a non-const reference to a private member (\texttt{int\& getVal()}) breaks encapsulation. Fix it.
\item \textbf{Lifetime.} Trace construction/destruction order for \texttt{\{ Student a(1,"A",4); Student b=a; \}} How many constructor/destructor calls? Which kind?
\end{enumerate}
% End of Chapter 5
